\scanpage{196}
\[
\underbrace{\varphi_{n'}\underbrace{\mathcal A_*(\tau)\mathcal A_*(\sigma)}_{\mathcal A_*(\sigma\tau)}}_{\mathcal A_*(\sigma\tau)\varphi\ (\text{par hyp. sur }\varphi_{n'})}
=
\underbrace{\mathcal A_*(\tau)\underbrace{\varphi_n\mathcal A_*(\sigma)}_{\mathcal A_*(\sigma)\varphi\ (\text{par hyp. sur }\varphi_n)}}_{\mathcal A_*(\sigma\tau)\varphi}
\qquad\forall\sigma:\Delta_n\to\Delta_1
\]
OK.

\uline{Dém. de la prop} Le cas $n=0$ est \unsure{triv.}, pour le cas $n=1$ il faut considérer les trois cas $\sigma=\mathrm{id}$ qui donne $\varphi_1=\varphi$, puis $\sigma$ application constante de valeur $0$ ($\alpha_0$) ou $1$ ($\beta_0$), de sorte que $\mathcal A_*(\alpha_0)=\alpha$, $\mathcal A_*(\beta_0)=\beta$, et le fait que $\varphi_1=\varphi$ y soit \unsure{compatible} n'est autre que l'hyp. suppl. sur $\varphi$ sur les augmentations \inserted{$\alpha,\beta$}. On peut donc supposer $n\ge2$. Considérons les applications \inserted{croiss} $\sigma:\Delta_n\to\Delta_1$; elles sont surj. ou constantes. Considérons les surjectives, il y en a $n$, soit $\sigma_{\unsure{r}}$ ($0\le r\le n-1$)
\[
\sigma_r(0)=0,\ -\ \sigma_r(r)=0,\ \sigma_r(r+1)=1,\ -\ \sigma_r(n)=1
\]
donc \deleted{$\unsure{\mathcal A}$}
\[
\begin{cases}
\mathcal A_*(\sigma_r)(X_0)=X_0+-+X_r\overset{\mathrm{dfn}}=U_r\\
\mathcal A_*(\sigma_r)(X_1)=X_{r+1}\ -+X_n
\end{cases}
\]
On doit donc avoir
\[
\mathcal A_*(\sigma_r)(X_0^{(i)})=U_r^{(i)}
\]
\uline{Lemme} Les $U_r^{(i)}$ (\deleted{$\unsure{i}$} $i\ge1$, $0\le r\le n-1$) engendrent la \inserted{$S$-}algèbre $\mathcal A_n^0$, qui est isomorphe en fait à $\unsure{S}\{U_0,-,U_{n-1}\}$,
\margin{195}
